\documentclass[10pt,a4paper]{amsart}
\usepackage{amsmath,amssymb,amsthm,graphicx}
\RequirePackage{amssymb,amsmath,amsthm,bm}
\RequirePackage{thmtools}
\RequirePackage[capitalize]{cleveref}
\providecommand{\tightlist}{%
  \setlength{\itemsep}{0pt}\setlength{\parskip}{0pt}}

%%%%% varphi instead of phi
\renewcommand{\phi}{\varphi}

\mathchardef\mhyph="2D 
\providecommand{\A}{\mathcal{A}}
\providecommand{\C}{\mathcal{C}}
\providecommand{\L}{\mathcal{L}}
\providecommand{\B}{\mathcal{B}}
\providecommand{\M}{\mathcal{M}}
\renewcommand{\S}{\mathcal{S}}
\providecommand{\ol}[1]{\overline{#1}}
\providecommand{\mbf}[1]{\mathbf{#1}}
\providecommand{\mf}[1]{\mathfrak{#1}}
\providecommand{\mrm}[1]{\mathrm{#1}}
\providecommand{\mc}[1]{\mathcal{#1}}
\providecommand{\ra}{\rightarrow}
\providecommand{\Ra}{\Rightarrow}
\providecommand{\La}{\Leftarrow}
\providecommand{\lr}{\leftrightarrow}
\providecommand{\hf}{\mathbb{HF}}
\providecommand{\spec}{DgSp}
\providecommand{\bSigma}{\bm\Sigma}
\providecommand{\bDelta}{\bm\Delta}
\providecommand{\bPi}{\bm\Pi}
\providecommand{\Pinf}[1]{\Pi^{\mathrm{in}}_{#1}}
\providecommand{\Sinf}[1]{\Sigma^{\mathrm{in}}_{#1}}
\providecommand{\dSinf}[1]{d\mhyph\Sigma^{\mathrm{in}}_{#1}}
\providecommand{\Dinf}[1]{\Delta^{\mathrm{in}}_{#1}}
\providecommand{\bp}{\bar p}
\providecommand{\bq}{\bar q}
\providecommand{\ba}{\bar a}
\providecommand{\bb}{\bar b}
\providecommand{\bc}{\bar c}
\providecommand{\bd}{\bar d}
\providecommand{\bx}{\bar x}
\providecommand{\by}{\bar y}
\providecommand{\bz}{\bar z}
\providecommand{\LO}{\mathrm{LO}}
\declaretheorem[name={Theorem},numberwithin=section]{theorem}
\declaretheorem[name={Proposition}, sibling=theorem]{proposition}
\declaretheorem[name={Lemma}, sibling=theorem]{lemma}
\declaretheorem[name={Fact}, sibling=theorem]{fact}
\declaretheorem[name={Claim}, numberwithin=theorem]{claim}
\declaretheorem[name={Definition},numberwithin=section,sibling=theorem, style=definition]{definition}
\declaretheorem[name={Remark}, numbered=no,style=definition]{remark}
\declaretheorem[name={Corollary}, sibling=theorem]{corollary}
\declaretheorem[name={Question},sibling=theorem]{question}
\newcommand{\T}{{\mathcal{T}}}
\newcommand{\SR}{SR}
\newcommand{\pSR}{pSR}
\newcommand{\Q}{\mathbb{Q}}
\newcommand{\Z}{\mathbb{Z}}
\newcommand{\K}{\mathcal{K}}
\newcommand{\N}{\mathcal{N}}
\renewcommand{\L}{\mathcal{L}}
\providecommand{\Ginf}[1]{\Gamma^{\mathrm{in}}_{#1}}
\title{Multiplying by Z to the alpha preserves back-and-forth relations with a doubled shift}
\author{David Gonzalez, Dino Rossegger}
\date{arxiv:2305.07126v1; CC BY 4.0}
\begin{document}
\maketitle
\noindent\textit{Editorial note.} The counted statement is the author environment at source0.tex lines 1275-1278 and its proof environment opens at line 1280, the line immediately after; the pairing is by adjacency and the retained proof range 1280-1301 is the whole environment.\par\medskip
\noindent\textit{Editorial note.} The retained preamble is the author style file dino.sty (package lines, the theorem declaration support, the notation and tuple macros and the theorem environments) followed by the ordinal and rank macros of the manuscript itself, reproduced verbatim; this is the same recipe as the previously published card of this article, so no package outside it is requested and extra\_packages stays empty.\par\medskip
\noindent\textit{Editorial note.} Metadata provenance: source\_id and source\_version are the values recorded in this source licence evidence file (sources/scott-linear-orders-2305.07126v1/LICENSE\_EVIDENCE.json: arxiv:2305.07126v1 / 2305.07126v1), which the card points at directly; the title and author list are transcribed from the authors\_index and title keys of the source record of the pinned copy. The card carries no counter setting and imposes no original numbering, since this article has no label-number dump and its numbering cannot be reconstructed from the retained ranges alone.\par\medskip
\noindent\textit{Editorial note.} No context span and no external reference are needed: the retained statement and proof use no label or citation from outside their own ranges.\par\medskip
\noindent\textit{Editorial scope.} One original result of the pinned author manuscript, counted once: the rank shift by twice alpha for back-and-forth comparisons of products with Z to the alpha (author label lem:zpowersbfpreserve).\par\medskip
\noindent\textit{Allowed standard prerequisites (editorial).}\begin{itemize}
\item Multiplication of linear orders by Z to the alpha. Z\textasciicircum{}alpha times L is the lexicographic product of the order Z\textasciicircum{}alpha with L, whose elements are pairs (delta, l).. It is the operation the lemma relates to the back-and-forth relations.
\item Back-and-forth relations and their intervals. The rank beta relation between two orders is decided by the rank beta minus one relations between the intervals determined by the chosen tuples.. It is how the successor step is verified.
\end{itemize}\medskip
\begin{lemma}\label{lem:zpowersbfpreserve}
For any $\L$ and $\K$
$$\L\leq_\beta \K\implies \Z^\alpha\cdot \L\leq_{2\alpha+\beta} \Z^\alpha\cdot \K.$$
\end{lemma}

\begin{proof}
 We demonstrate this by induction on $\beta$. Notice that the base case is given by the above observation. 
 The limit case follows immediately by the definition of the back-and-forth relations along with the observation that for any linear order $\N$, $\L\leq_\beta \K \implies \N\cdot \L\leq_\beta \N\cdot \K$.
 Thus, we need only look at the successor case.

 Let $\beta=\gamma+1$ and assume $\L\leq_{\gamma+1} \K$. We need to show that
 $\Z^\alpha\cdot \L\leq_{2\alpha+\gamma+1}\Z^\alpha\cdot \K.$ We view
 $\Z^{\alpha}\cdot K$ as a product ordering with elements of the the form
 $(\delta,b)$ where $b\in K$ and $\delta\in \Z^{\alpha}$. Consider a play of the
 game where the $\forall$-player plays a tuple $(\delta_i,b_i)_{i\in k}$ from $\Z^\alpha\cdot
 \K$. If $s$ is a winning strategy for demonstrating that $\L\leq_{\gamma+1}
 \K$, the $\exists$-player may play $(\delta_i,s(b_i))_{i\in k}$ in response. Let
 $b_{-1}=s(b_{-1})=-\infty$ and $b_k=s(b_k)=\infty$ by convention. We need only
 show that, $$[(\delta_i,b_i),(\delta_{i+1},b_{i+1})] \leq_{2\alpha+\gamma}
 [(\delta_i,s(b_i)),(\delta_{i+1},s(b_{i+1}))].$$ Note that the left hand side
 is isomorphic to $\zeta_\alpha+\Z^\alpha\cdot(b_i,b_{i+1})+\zeta_\alpha^*$ and
 the right hand side is isomorphic to
 $\zeta_\alpha+\Z^\alpha\cdot(s(b_i),s(b_{i+1}))+\zeta_\alpha^*$. As $s$ is a
 winning strategy, by induction,
 \[\Z^\alpha\cdot(b_i,b_{i+1})\leq_{2\alpha+\gamma}
 \Z^\alpha\cdot(s(b_i),s(b_{i+1})),\] which demonstrates the claim.
\end{proof}

\end{document}
